\begin{formbox}
\textbf{Individual Contribution Statement --- Project 2 (PDE BVP)}

\textbf{Name:} Meichen Ding (2362344)

\textbf{Role this project:} Mathematical theory lead (consistency, adaptivity
and convergence): owner of the model formulation, the discrete scheme, its
truncation-error and interface analysis, and the discrete-structure results.

\textbf{The specific things I did this project:}
\begin{enumerate}[leftmargin=1.4em,itemsep=0.15em]
  \item
  Wrote \S1.2 and \S2.1--\S2.4 of the report's \path{Method.tex}: the continuous
  problem and its hypotheses, the flux-balance derivation that fixes
  one shared face coefficient per edge, the index map and block-tridiagonal
  structure of the assembled operator (with the non-zero count $5n^2-4n$), the
  symmetric Dirichlet elimination and the hand-computed $4\times4$ matrix at
  $n=2$ used as the small-case oracle, the layered exact solution, the
  second-order local truncation error with its $\theta$-correction, the
  series-resistance proposition for the interface conductance with the relative
  error $(1-c)^2/(4c)$, the SPD / $M$-matrix / discrete-maximum-principle
  theorem, and the discrete stability and $\mathcal{O}(h^2)$ convergence theorem
  with its maximum-principle barrier argument.
  \item
  Wrote the independent verification scripts \path{support/verify/}:
  \texttt{check\_interface\_order.py} (one-dimensional conservative scheme,
  $a_L=1$, $a_R=10$, $N=31,63,127,255$; harmonic nodes exact to $10^{-14}$ and
  every face flux equal to $-20/11$, arithmetic nodal error first order with
  measured orders $1.015/1.008/1.004$ and the closed-form chain flux
  reproduced), \texttt{check\_stability.py} (the entrywise-ratio counterexample
  at $a=1+9\sin^2(\pi x)$, ratio $3.10$; $\|A_h^{-1}\|_\infty$ against
  $1/(8a_{\min})$ for five coefficients) and
  \texttt{check\_matrix\_evidence.py} (symmetry defect, \texttt{eigsh} smallest
  eigenvalue and dense Cholesky on the submitted assembler). These produced the
  numbers in the appendix tables and in \S2.4.
  \item
  Ran the AI-assisted claim audits of Part A: the three-layer separation of the
  arithmetic-averaging errors (audit 1), the falsification of the
  Loewner-implies-entrywise argument and its replacement by the barrier proof
  (audit 2), and the separation of ``symmetric and conservative'' from
  ``physically correct'' (audit 3), including the numerical values that decide
  each one ($0.662$ versus $3.85\times10^{-13}$; $3.10>1$; $0.0736$ versus
  $0.125$).
\end{enumerate}

\textbf{What I would do differently:}
I would pin down the mesh convention with the code owner before writing any
convergence sentence. The skeleton uses $h=1/(n+1)$ with $n$ interior points
while the report uses $h=1/N$ with $N$ intervals, so the same symbol $N$ means
different things in the skeleton, in the layered tables and in the FEM
comparison; the interface-parity requirement $N$ odd for the FDM and $N$ even
for the FEM then has to be stated as design rather than discovered late. I would
also have demanded the two-dimensional layered numbers before writing the
one-dimensional flux estimate, because that is exactly where the two-layer
statement became necessary.

\textbf{Anything I should have asked for help with earlier:}
I should have asked, in Week 5, for each member to commit their own sections and
their own ICS from their own account. By the time the appendices were written
the repository history was a single upload commit, so authorship had to be
reconstructed from memory and from section ownership instead of being visible in
\texttt{git log}. I also should have asked for a review of the stability proof
before it went into a draft: the Loewner-order step survived an internal
re-reading and was only caught when it was written out as a counterexample
search.

\vspace{0.8em}
\textbf{Signature:} \underline{\hspace{5cm}} \hfill
\textbf{Date:} \underline{\hspace{3cm}}
\end{formbox}
